% Draft (2026-11-07). Numbers are from preregistered studies already run (Appendix A and E3, Amendment 3).
\section{Transfer: scientific code, results and data}
\label{sec:transfer}

Section~\ref{sec:framework} separates capacity misses, where the evidence decides the case but the auditor does
not use it, from information misses, where the evidence does not decide the case. Three preregistered studies of
scientific work fall on either side of this distinction. They preceded the experiment above and are reported in
full in Appendix~\ref{app:earlier}. Each uses ground truth that no model supplies.

\paragraph{Code: what the auditor is told it is reviewing.} SciCode \citep{tian2024scicode} splits research problems
into steps with tests written by scientists. We label generated steps that pass the tests correct and those that fail defective (the tests can themselves
reject correct code; \citealp{hu2026scicodeverified}) (81 defective, 150 correct instances). The auditor sees the step's text and code but no tests.
Given the whole problem as the task, a cross-vendor auditor at eight readings flagged 87.7\% [72.7, 98.6] of
defective and 68.7\% [56.6, 79.7] of correct steps. Some objections concerned a single step's failure to implement
the whole problem. A second registered arm changed only the task text to name the step as the deliverable;
the program bytes were identical. Flags on correct code fell by 32.7 points [22.4, 42.8]. As on the constructed
units, telling the auditor what was required produced the largest change, rather than reading more.

\paragraph{Results: where the evidence decides, computation helps.} For 74 correct SciCode programs we built a run
log, a results file and a report in which every number cites its source. We injected either a reporting fault
(the artefacts disagree) or a fabrication written consistently into all of them. A deterministic provenance
check flagged all 49 reporting faults and none of the 25 fabrications, whose artefacts all agreed.
The model auditor flagged 23 of the 25, usually by recomputing the result from the code and inputs (exact McNemar
$p=2.4\times10^{-7}$). Re-execution uses the same information and flagged 70 of the 74 faulty items.
It missed four whose change fell inside the default numerical tolerance. The evidence here determines the
right value, so a miss is a capacity miss. More capacity, through a model that computes or simply running the
code, removes most of them.

\paragraph{Data: where the documentation carries the information.} We gave each of twelve public regression tables a
data card. Each table was either left clean or given one injected fault (mixed units, negated values, duplicated
rows, a shuffled target, swapped columns, sentinel values, rounding). We held out eight tables (266 items) and
combined a validator written mechanically from the card with a model auditor (GPT or Claude, with a shell on the
full table or reading 120 rows). Adding the model raised correct localisation, scored by code against the injection manifest, by 2.3 points
[0.4, 4.8] and flags on clean items by 16.0 points [3.2, 31.4]. Neither the full table nor tool use helped.
In an unsandboxed pilot, one auditor searched the disk and the web for the original files. Every subsequent
reading ran in an operating-system sandbox. Even there, 57\% of that auditor's held-out readings attempted
out-of-scope access; these attempts were logged and returned nothing. An earlier study on four tables, at a clean-flag budget fixed in advance, found
a larger gain in findings naming the faulty column (10.7 points [6.0, 15.9]) over a simpler validator.
In both studies, where the documentation already states what valid data look like, a check written from it
finds most of the faults that it and the model find together.
